\documentclass[14pt, a4paper]{extarticle} % ใช้ extarticle เพื่อรองรับ 14pt

% --- Packages สำหรับภาษาไทย (XeLaTeX) ---
\usepackage{fontspec}
\usepackage{polyglossia}
\usepackage{setspace} % สำหรับปรับระยะห่างระหว่างบรรทัด

\setmainlanguage{thai}
\setotherlanguage{english}
\setstretch{1.25} % เพิ่มระยะบรรทัดให้สระบน-ล่างไม่เบียดกัน

% การตั้งค่าฟอนต์สารบรรณพร้อมขยายขนาด (Scale)
\defaultfontfeatures{Scale=1.25} 
\newfontfamily\thaifont{TH Sarabun New}[
  UprightFont = *,
  BoldFont = * Bold,
  ItalicFont = * Italic,
  BoldItalicFont = * Bold Italic
]
\newfontfamily\thaifontsf{TH Sarabun New}[UprightFont = *]
\newfontfamily\thaifonttt{TH Sarabun New}[UprightFont = *]

\setmainfont{TH Sarabun New}
\setsansfont{TH Sarabun New}
\setmonofont{TH Sarabun New}

% --- Packages เสริม ---
\usepackage{geometry}
\usepackage{graphicx}
\usepackage{hyperref}
\usepackage{listings}
\usepackage{xcolor}
\usepackage{titlesec}

% ปรับ Margin ให้สมดุลกับตัวอักษรที่ใหญ่ขึ้น
\geometry{top=0.8in, bottom=0.8in, left=0.8in, right=0.8in}
\hypersetup{colorlinks=true, linkcolor=blue, urlcolor=cyan}

% ปรับขนาดหัวข้อให้เด่นชัด
\titleformat{\section}{\normalfont\Large\bfseries}{\thesection}{1em}{}
\titleformat{\subsection}{\normalfont\large\bfseries}{\thesubsection}{1em}{}

\title{\textbf{คู่มือระบบวิเคราะห์ผลการบิน RASAT/FAI \\ (Paramotor Triangle Scoring)}}
\author{Competition Director Manual}
\date{\today}

\begin{document}

\maketitle
\newpage
\tableofcontents
\newpage

\section{บทนำและวัตถุประสงค์}
ระบบนี้พัฒนาขึ้นเพื่อเพิ่มความแม่นยำและความโปร่งใสในการตัดสินการแข่งขัน Paramotor ประเภท Triangle โดยใช้ระบบการประมวลผลแบบ Batch เพื่อลดภาระของกรรมการในการวิเคราะห์ไฟล์ IGC ทีละคน และเพื่อให้มั่นใจว่านักกีฬาจะได้รับระยะทางที่ยุติธรรมที่สุดตามกติกา FAI

\section{แนวคิดหลักและการคิดคะแนน (Analysis Concepts)}

\subsection{กติกา FAI Triangle 28\%}
หัวใจสำคัญของกติกาคือสัดส่วนของสามเหลี่ยม:
\begin{itemize}
    \item \textbf{FAI Triangle (1.5x)}: ด้านที่สั้นที่สุดต้องมีระยะทางอย่างน้อย 28\% ของระยะทางรวม 3 ด้าน
    \item \textbf{Flat Triangle (1.0x)}: หากไม่ผ่านเกณฑ์ 28\% จะถือเป็นสามเหลี่ยมทั่วไปและใช้ตัวคูณปกติ
\end{itemize}



\subsection{การตรวจสอบ Hidden Gates (CTE)}
ระบบจะวางจุดตรวจสอบ (Gates) ไว้บนเส้นทางสามเหลี่ยมในทุก ๆ 1 กม. หากนักบินบินผ่านในรัศมีที่กำหนด (ปกติ 200 ม.) จะถือว่า \textbf{Scored}

\section{ขั้นตอนการใช้งานสำหรับ CD}

\subsection{การตั้งค่า Task (Configuration)}
ก่อนเริ่มแข่ง CD ต้องแก้ไขพิกัดในไฟล์ \texttt{task\_config.json}:
\begin{itemize}
    \item \textbf{start\_point}: พิกัดจุดเริ่ม (Latitude, Longitude)
    \item \textbf{finish\_point}: พิกัดจุดจบ (หากเป็น Closed Triangle ให้ใส่เหมือนจุดเริ่ม)
    \item \textbf{sp\_radius\_meters}: รัศมีจุดเริ่ม/จบ (ปกติ 200 ม.)
\end{itemize}

\subsection{การประมวลผลไฟล์ IGC}
\begin{enumerate}
    \item รันสคริปต์ \texttt{setup.bat} (Windows) หรือ \texttt{setup.sh} (Mac/Linux)
    \item เมื่อหน้า GUI ปรากฏ ให้กดปุ่ม \textbf{Select IGC Files}
    \item เลือกไฟล์ IGC ของนักกีฬาทุกคนพร้อมกัน ระบบจะทำการ Copy ไฟล์ไปยังโฟลเดอร์ \texttt{igcFiles} อัตโนมัติ
\end{enumerate}

\section{การตีความผลลัพธ์ (Interpretation)}

\subsection{ตารางสรุปคะแนน (CSV Result)}
เปิดไฟล์ \texttt{competition\_results.csv} เพื่อดูอันดับ:
\begin{itemize}
    \item \textbf{Effective\_km}: ระยะทางสุทธิที่ใช้จัดลำดับ (รวมตัวคูณแล้ว)
    \item \textbf{Status}: หากขึ้น \texttt{FAILED} แสดงว่าบินไม่เข้าจุดเริ่มหรือจุดจบ
\end{itemize}

\subsection{ภาพวิเคราะห์เส้นทาง (Visualization)}
เปิดดูรูปในโฟลเดอร์ \texttt{results\_track\_analysis}:
\begin{itemize}
    \item \textbf{จุดสีเขียว}: ผ่าน Hidden Gate สำเร็จ
    \item \textbf{จุดสีแดง}: พลาด Hidden Gate (บินหลุดแนวเส้นทาง)
\end{itemize}



\end{document}